\chapter{化学反应速度和化学平衡}
我们已经学习了一些化学反应，知道化学反应往往需要在一定的条件下进行，例如，上章讲到使氢气和氮气化合生成氨时，就要在高温、高压和有催化剂存在的条件下进行。
为什么一个反应的进行需要这样或那样的条件呢？
这就要从以下两个方面来认识: 一个是反应进行的快慢，即化学反应速度问题; 一个是反应进行的程度，也就是达到化学平衡的问题。
这两个问题不仅是今后学习化学所必需的基础理论知识，也是研究化工生产适宜条件时需要掌握的化学变化的规律。

\section{化学反应速度}
\subsection{化学反应速度}
各种化学反应的进行有快有慢，如氢、氧爆鸣气的爆炸反应、酸碱溶液的中和反应等瞬时就能完成; 而石油的形成就要经过亿万年。
这里所说的快和慢，是指在一定条件下进行的化学反应。
只要掌握了化学反应的规律，我们就可以根据科学研究、生产和生活的需要，采取适当的措施，把原来进行得很慢的反应加快，如加速炼钢过程，加速合成树脂或合成橡胶的反应等; 或把原来进行得快的反应减慢，如钢铁的防锈，塑料、橡胶的防止老化等等。
这些对发展生产，节约支出，促进社会主义现代化建设，都具有重要的意义。

怎样来衡量化学反应速度的大小呢？

通常对于某些反应可以通过观察反应物的消失速度和生成物的出现速度作出对这个反应速度的定性判断。
例如，当把一条镁带放在一个盛稀盐酸的烧杯里，在镁带（反应物）迅速消失的同时，会很快地放出氢气（生成物）。
当在同浓度的盐酸中放入铁块时，铁块以较慢速度消失，同时氢气放出的速度也较缓慢。
显然第一个反应的速度比第二个反应大得多。

怎样定量地较准确地表示化学反应速度呢？
化学反应的速度是用单位时间（如每秒、每分或每小时等）内反应物或生成物的量（\unit{mol}）的变化来表示，通常是用单位时间内反应物浓度的减小或生成物浓度的增大来表示。
浓度的单位一般为 \unit{mol/L}，反应速度的单位就是 \unit{mol.L^{-1}.min^{-1}} 或 \unit{mol/L^{-1}.s^{-1}} 等。

例如: 某一反应物的浓度是 \qty{2}{mol/L}，经过两分钟的反应后，它的浓度变成了 \qty{1.8}{mol/L}，即两分钟后反应物的浓度减小了 \qty{0.2}{mol/L}，这就是说，在这两分钟内它的平均反应速度为 \qty{0.1}{mol.L^{-1}.min^{-1}}。

\subsection{影响反应速度的条件}
不同的化学反应，具有不同的反应速度。
例如，酸和碱的中和反应，就比氮分子跟氢分子合成氨的反应快得多，这说明参加反应物质的性质是决定化学反应速度的主要因素。
但是外界条件对化学反应速度也有一定的影响，就是同一个化学反应，如浓度、压强、温度、催化剂等外界条件不同时，反应速度也不相同。
下面我们来研究影响化学反应速度的几个重要条件。

\subsubsection{浓度对化学反应速度的影响}
我们在初中学习氧的性质时，曾看到硫在空气中缓慢燃烧并产生微弱的淡蓝色火焰，在纯氧中迅速燃烧并发出明亮的蓝紫色火焰。
这说明硫在纯氧中跟氧化合的反应比在空气中的反应进行得更快、更剧烈。
这是因为纯氧中氧分子的浓度比空气中氧分子的浓度大的缘故。
从下面的实验可以看出，在溶液中进行的反应也是这样的情况。

\begin{Experiment}
  取两个试管，在第一个试管中加入 \qty{0.1}{M} \ce{Na2S2O3}（硫代硫酸钠）溶液 \qty{10}{ml}；在第二个试管中加入 \qty{0.1}{M} \ce{Na2S2O3} 溶液和蒸馏水各 \qty{5}{ml}。

  另取两个试管，每个试管加入 \qty{0.1}{M} \ce{H2SO4} 溶液各 \qty{10}{ml}。然后同时分别倒入上面两个盛有 \ce{Na2S2O3} 溶液的试管中，并注意观察两个盛有混和溶液试管里出现浑浊现象的先后。
\end{Experiment}

实验结果表明: 首先出现浑浊现象的是第一个试管，接着才是第二个试管。

通过许多实验证明，当其它条件不变时，增加反应物的浓度，可以增大反应的速度。

\subsubsection{压强对化学反应速度的影响}
对于气体来说，当温度一定时，一定量气体的体积与其所受的压强成反比。
这就是说，如果气体的压强增大到原来的两倍，气体的体积就缩小到原来的一半，单位体积内的分子数就增多到原来的两倍，如\cref{fig:3-1} 所示。
所以，增大压强，就是增加单位体积里反应物的摩尔数，即是增大反应物的浓度，因而可以增大反应的速度。
相反，减小压强，气体的体积就扩大，浓度减小，因而反应速度减小。
\begin{figure}
  \includegraphics{3-1.pdf}
  \caption{压强大小与一定量分子所占体积的示意图}\label{fig:3-1}
\end{figure}

如果参加反应的物质是固体、液体或溶液时，由于改变压强对它们的体积改变很小，因而对它们的浓度改变很小，可以认为压强与它们的反应速度无关。

\subsubsection{温度对化学反应速度的影响\quad 活化能}
我们知道，许多化学反应都是在加热的情况下发生的。
例如，在常温下煤在空气里甚至在纯氧里不能燃烧，只有加热到一定温度时才能燃烧，并且越燃越旺。
物质在溶液里进行的反应也有类似的情况。
\begin{Experiment}
取两个试管，在每个试管中各加入 \qty{0.1}{M} \ce{Na2S2O3} 溶液 \qty{10}{ml}。另取两个试管，在每个试管里各加入 \qty{0.1}{M} \ce{H2SO4} 溶液 \qty{10}{ml}。取一个盛有 \ce{Na2S2O3} 溶液的试管和另一个盛有 \ce{H2SO4} 溶液的试管组成一组，即四个试管组成两组。

然后将一组试管插入冷水里，另一组试管插入热水里。过一会儿，同时分别将两组试管里的溶液混和，并仔细观察热水和冷水中盛混和溶液的试管里出现浑浊现象的情况。
\end{Experiment}

实验结果表明：插在热水中盛混和溶液的试管里首先出现浑浊的现象，接着才是插在冷水中盛混和溶液的试管里出现浑浊。
这是什么原因呢？因为 \ce{Na2S2O3} 溶液跟稀 \ce{H2SO4} 作用时发生如下的反应：
\[ \ce{ Na2S2O3 + H2SO4 \xlongequal{\quad} Na2SO4 + SO2 + S v + H2O } \]
或
\[ \ce{ S2O3^2- +2H+ \xlongequal{\quad} SO2 + S v +H2O } \]
反应的快慢可借反应生成硫（不溶于水，使溶液浑浊）所需时间的长短来量度。
插入热水中的试管由于温度高，反应快，所以先出现浑浊现象; 插入冷水中的试管由于温度低，反应慢，后出现浑浊现象。
由此可见，温度升高，化学反应一般要加快。
经过多次实验测得，温度每升高 \qty{10}{\celsius}，反应速度通常增大到原来的 \numrange{2}{4} 倍。

为什么改变反应物的浓度或改变反应的温度时，就会改变化学反应的速度呢？

发生化学反应的先决条件是反应物的分子（或离子）必须互相接触，互相碰撞，否则，就不可能发生化学反应。
以气体的反应为例，任何气体中分子间的碰撞数都是非常巨大的。
在 \qty{1}{atm} 和 \qty{500}{\celsius} 时，浓度为 \qty{0.001}{mol/L} 的 \ce{HI} 气体，分子碰撞次数每升每秒达 \num{3.5e28} 之多，如果每次碰撞都能发生化学反应，\ce{HI} 分解的反应瞬间就可完成，但事实并不是这样的。
又如，在常温常压下，氢气和氧气的混和物可以长时间放置而不发生明显的反应。
可见反应物的分子的每次碰撞不一定都能发生化学反应，能够发生化学反应的碰撞是很少的，我们把这种能够发生化学反应的碰撞叫做\Concept{有效碰撞}。

为什么分子碰撞时有的能发生化学反应，有的不能呢？因为化学反应的过程，就是反应物分子中的原子重新组合而成生成物分子的过程，也是反应物分子中化学键的断裂，生成物分子中化学键的形成的过程。
在这过程中，反应物分子必须具有足够的能量，碰撞时才能使化学键削弱或断裂，从而发生化学反应。
可见发生有效碰撞的分子必须具有较高的能量。 
在一定温度下，气体分子具有一定的平均能量，但并不是所有的分子都具有相等的能量，实际上有的分子的能量比平均能量高，有的比平均能量低。
对某一反应来说，必有一部分具有比平均能量高的分子，它们才能发生有效碰撞。
这种分子就是\Concept{活化分子}。
活化分子具有的最低能量与分子平均能量的差，叫做\Concept{活化能}。
如\cref{fig:3-2} 所示：
\begin{figure}
  \includegraphics{3-2.pdf}
  \caption{反应过程中能量变化示意图}\label{fig:3-2}
\end{figure}
具有平均能量的反应物分子，要吸收 $E_1$ 的能量，才能变成活化分子，这个能量 $E_1$ 就是该反应的活化能。
活化分子变成生成物分子的过程中，要放出 $E_2$ 的能量。
显然图中 $E_2>E_1$，在反应物变成生成物的整个过程中要放出 $E_2-E_1$ 的能量，这个能量就是反应热。
\cref{fig:3-2} 所示的反应是放热反应。
无论放热反应或吸热反应，都需要有一个活化的过程。
前面提到的在常温下煤在空气里必须加热才能燃烧，就是因为加热能促使煤和氧气中的分子活化，发生氧化反应，反应发生后，就可由放出的反应热维持反应继续进行，而不需要外部再供给能量。

不同的化学反应所需的活化能不同。
如果反应的活化能较低，则在一定温度下，活化分子百分数就较大，反应就比较容易进行。

在其它条件不变时，对某一反应来说，反应物分子中活化分子的百分数是一定的，因此，单位体积内活化分子的数目与单位体积内反应物分子的总数成正比，也就是和反应物的浓度成正比。
当反应物浓度增大时，单位体积内分子增多，活化分子数也相应增大。
譬如原来 1 单位体积里有 100 个反应物的分子，其中只有 5 个活化分子，如果每 1 单位体积内的反应物分子增加到 200 个，其中必定有 10 个活化分子，那么单位时间内的有效碰撞次数也相应增多，反应就快。

在浓度一定时，升高温度，反应物的分子的能量增加，必然有一部分原来能量较低的分子变成了活化分子，从而增加了反应物分子中活化分子的百分数，有效碰撞次数增多了，因而加大了反应的速度。
当然，由于温度升高，会使分子的运动加快，这样单位时间内反应物分子间的碰撞次数增加，反应也会相应地加快，但这不是反应加快的主要原因。

\subsubsection{催化剂对反应速度的影响}
关于催化剂和催化作用的初步知识，在氧气的实验室制法、硫酸的工业制法和氨的催化氧化等教材里都学习过了。
现在我们要学习催化剂与化学反应速度的关系。
\begin{Experiment}
在两个试管里分别加入 3\% 的过氧化氢（\ce{H2O2}）溶液 \qty{3}{ml} 和合成洗涤剂（产生泡沫以示有气体生成）溶液 3 或 4 滴。
在其中的一个试管里加入少量 \ce{MnO2}，观察两个试管里的反应现象有什么不同。
\end{Experiment}
实验结果表明: 在放有少量 \ce{MnO2} 的试管中，很快有气泡生成,而没有放 \ce{MnO2} 的试管里气泡产生得慢而且少。
这是由于 \ce{MnO2} 能加快 \ce{H2O2} 的分解; 使氧气产生得快而且多，反应后 \ce{MnO2} 的组成和质量都没有变化。
可见 \ce{MnO2} 在 \ce{H2O2} 的分解反应中起了催化作用。
\[ \ce{ 2H2O2 \xlongequal{\ce{MnO2}} 2H2O + O2 ^ }\]

催化剂能够增大反应速度的原因，是它能够降低反应的活化能，从而使更多的反应物分子成为活化分子，大大增加单位体积内反应物分子中活化分子的百分数，成千成万倍地增大化学反应的速度。
从\cref{fig:3-3} 可以清楚地看出催化剂与反应速度的关系。
图中实线曲线表示不用催化剂时，反应物变成生成物的过程，$E_1$ 是不用催化剂时的反应的活化能; 虚线曲线表示用催化剂时,反应物变成生成物的过程，$E_2$ 是有催化剂存在时该反应的活化能。
显然，使用催化剂能使反应的活化能降低，因而能大大增大反应的速度。
这好象人们要翻越一座陡峭的高山，很费劲、很慢，能翻越的人就少些; 如果在山下开凿一条隧道，穿过隧道就容易、就快，能穿过的人就多一样。

\begin{figure}
  \includegraphics{3-3.pdf}
  \caption{催化剂与活化能的关系示意图}\label{fig:3-3}
\end{figure}

\begin{Reading}[补充阅读]{}
为什么使用适当的催化剂能够降低反应的活化能呢？ 现以接触法制造硫酸使用的接触剂 \ce{V2O5}，和合成氨工艺中使用的铁触媒为例，来说明它们是怎样降低活化能的。

制造硫酸的过程中，关键是怎样把 \ce{SO2} 更快地转变为 \ce{SO3}。如果使用 \ce{V2O5} 做催化剂时，有的化学工作者认为反应是按以下两步进行的：
\begin{gather*}
  \ce{ SO2 + V2O5 \xlongequal{\quad} SO3 + V2O4 }\\
  \ce{ V2O4 + SO2 + O2 \xlongequal{\quad} SO3 + V2O5 }
\end{gather*}

可见使用 \ce{V2O5} 做催化剂，能使一步进行的反应变成为分两步进行的反应。
实验测得这样将会使反应活化能降低到原来的二分之一左右，从而使反应速度大大加快。
由此可见，催化剂并不是不参与化学反应，只是反应前与反应后的组成和质量未改变罢了。

在氢跟氮合成氨的过程中，如果使用铁触媒时，铁触媒能吸附氢分子和氮分子，同时削弱氮分子和氢分子里的化学键，化学键被削弱了的这两种分子互相碰撞而合成氨分子。
由于使用铁触媒，合成氨反应的活化能降低到原来的四分之一左右。
同时由于吸附，在铁触媒表面上反应物浓度增大了，也能增大反应的速度。
这样，使用铁触媒的反应速度比无催化剂时成万倍地增大。
\end{Reading}

催化剂在现代化学和化工生产中占有极为重要的地位。 
据初步统计约有 85\% 的化学反应需要使用催化剂。
尤其在当前大型化工生产、石油化学工业生产中，很多反应还必须靠使用性能优良的催化剂来实现。
但是应当注意，因为有一些物质，即使是很少量的，当混入催化剂中时，就会急剧降低甚至破坏催化剂的催化能力，这种作用叫做催化剂的中毒。
为了防止催化剂的中毒，需要对原料进行一系列的净化过程。

影响化学反应速度的条件很多，除了温度、浓度、压强（有气体参加的反应）、催化剂以外，还有光、超声波、激光、放射线、电磁波、反应物颗粒的大小、扩散速度、溶剂等等，对某些化学反应来说，也是影响反应速度的一些重要条件。
例如，煤粉的燃烧就比煤块快得多，溴化银见光很快分解等等。

\begin{Theorem}{讨论}
举出你学过或生活中接触过的实例，说明影响反应速度的条件。
\end{Theorem}

以上研究了影响反应速度的一些重要条件，但在化学研究和化工生产中，只考虑反应速度是不够的。
例如，在合成氨工业中除了要求氮和氢尽可能快地转变为氨外，还要使氮和氢尽可能多地转变为氨。
这就涉及到化学反应的另一个问题一化学平衡。
下节就要研究关于化学平衡的问题。

\begin{Practice}[习题]
\begin{question}
  \item 举例说明加快和减慢反应速度的实际意义。
  \item 在下列反应
  \[ \ce{CO + H2O <=> CO2 + H2} \]
  里，起始浓度 $[\ce{CO}]=[\ce{H2O}]=\qty{0.02}{mol/L}$，\qty{1}{min} 后测得 $[\ce{CO}]=\qty{0.005}{mol/L}$，求分别以 \ce{CO} 和 \ce{H2} 的浓度表示的反应速度。
  \item 为什么反应的活化能低，反应速度就较大呢？
  \item 在空气中燃烧木炭是一个放热反应，为什么木炭燃烧时必须先引火点燃？点燃后停止加热，能够继续燃烧吗？为什么？
  \item 在一块大理石（主要成分是 \ce{CaCO3}）上,先后滴加 \qty{1}{M} \ce{HCl} 和 \qty{0.1}{M} \ce{HCl}，哪个反应快？先后滴加同浓度的热盐酸和冷盐酸，哪个反应快？用大理石块和大理石粉跟同浓度的盐酸起反应，哪个反应快？
  \item 在进行铁跟盐酸的反应时，用什么方法可以使反应加快？ 设计一个实验方法来粗略地衡量反应的速度。
  \item 为什么升高温度和增加反应物的浓度，都能增大反应的速度呢？
  \item 为什么使用适当的催化剂能够使一些反应的速度增大呢？
  \item 举例说明催化剂对化工生产的重要意义。
  \item 为什么增加反应物分子中活化分子的百分数能够增大反应的速度呢？用哪些方法可以增加反应物分子中活化分子的百分数呢？
\end{question}
\end{Practice}

\section{化学平衡}
现在我们来学习化学平衡，化学平衡主要是研究可逆反应的规律，如反应进行的程度以及各种条件对反应进行情况的影响等等。

\subsection{化学平衡是动态平衡}
固态溶质溶解在液态溶剂里，当形成饱和溶液时，达到溶解平衡。
这时，溶解的过程并没有停止，只是溶解和结晶的速度相等罢了。
因此，溶解平衡是一种动态平衡。
那么，在化学反应中，我们将要研究的可逆反应的情形又是怎样的呢？
我们可以通过下面的反应来进行研究。
\[ \ce{ 2SO2 + O2 <=>[\text{催化剂}][\triangle] 2SO3} \]

在 \qty{500}{\celsius}，\qty{1}{atm} 标准大气压时，把 2 体积的二氧化硫和 1 体积的氧气的混和物，通入一个装有催化剂的密闭容器里。
反应开始后，\ce{SO2} 和 \ce{O2} 的浓度逐渐减小，\ce{SO3} 的浓度逐渐增大，最后能得到含 91\%（体积组成）三氧化硫的混和气体。
这时候，容器里反应物 \ce{SO2}、\ce{O2} 和生成物 \ce{SO3} 在混和物中的浓度就不再发生变化。

为什么反应进行到一定程度，生成物的含量就不继续增大了呢？
也就是说，反应为什么不能进行到底呢？
我们可以用可逆反应中正反应和逆反应的反应速度的变化来说明。

当反应开始的时候，\ce{SO2} 和 \ce{O2} 的浓度最大，因而它们化合生成 \ce{SO3} 的正反应速度最大; 而 \ce{SO3} 的浓度为零，因而它分解生成 \ce{SO2} 和 \ce{O2} 的逆反应的速度也是零。
以后，随着反应的进行，反应物 \ce{SO2} 和 \ce{O2} 的浓度逐渐减小，正反应的速度就逐渐减小；生成物 \ce{SO3} 的浓度逐渐增大，逆反应的速度就逐渐增大。

如果外界条件不发生变化，化学反应进行到一定程度的时候，正反应和逆反应的速度相等，反应物和生成物的浓度不再发生变化。这时反应物和生成物的混和物（简称反应混和物）就处于\Concept{化学平衡状态}。

如果我们不是从 \ce{SO2} 和 \ce{O2} 的混和物开始反应，而是使纯净的 \ce{SO3}，在同样的温度（\qty{500}{\celsius}）、压强（\qty{1}{atm}）和催化剂存在的条件下起反应，\ce{SO3} 分解为 \ce{SO2} 和 \ce{O2}，当达到平衡状态时，混和气体里仍然含 91\%（体积）的 \ce{SO3}。

当反应达到平衡的时候，正反应和逆反应都仍在继续进行，只是由于在同一瞬间正反应 \ce{SO2} 和 \ce{O2} 化合生成的 \ce{SO3} 分子数和逆反应所分解的 \ce{SO3} 分子数相等，亦即正、逆反应的速度相等，所以反应混和物中各成分的百分含量不变。
因此，化学平衡是一种动态平衡。

\emph{化学平衡状态就是指在一定条件下的可逆反应里，正反应和逆反应的速度相等，反应混和物中各组成成分的百分含量保持不变的状态}。

\subsection{化学平衡常数}
下面我们以一氧化碳和水蒸气在高温时进行反应生成二氧化碳和氢气为例，来研究化学平衡的性质。

在密闭容器里这个反应不能进行到底，而是达到了化学平衡状态。
\[ \ce{ CO\,\text{（气）} + H2O\,\text{（气）}  <=>[\text{催化剂}][\text{高温}] CO2\,\text{（气）} + H2\,\text{（气）} }\]

如果把 \qty{0.01}{mol} 的 \ce{CO} 和 \qty{0.01}{mol} 的 \ce{H2O}（气）放在 \qty{1}{L} 的密闭容器里，加热到 \qty{800}{\celsius}，实验测得当生成物 \ce{CO2} 和 \ce{H2} 各为 \qty{0.005}{mol} 时，反应即达到了平衡状态。
达到了平衡的反应混和物（简称平衡混和物）里还有 \qty{0.005}{mol} 的 \ce{CO} 和 \qty{0.005}{mol} 的 \ce{H2O}（气）。这可以表示如下：

\noindent{\small
\begin{tblr}{colspec={rcX[c]X[c]c},hlines=0pt,vlines=0pt,
  vline{3}={1}{text=\clap{\ce{+}}},
  vline{4}={1}{text=\clap{\ce{<=>[\qty{800}{\celsius}]}}},
  vline{5}={1}{text=\clap{\ce{+}}},
  }
  & \ce{CO} & \ce{H2O}（气） & \ce{CO2}& \ce{H2} \\
  反应开始时，各物质的摩尔数 & 0.01  & 0.01  & 0     & 0     \\
  达到平衡时，各物质的摩尔数 & 0.005 & 0.005 & 0.005 & 0.005 \\
\end{tblr}
}

由于密闭容器的容积是 \qty{1}{L}，因此各物质的摩尔数在数值上跟它们的摩尔浓度（以摩尔/升为单位）相同。如果我们用符号 \ce{[\quad]} 表示某物质的浓度，那么上面的例子也可以写成：

\noindent{\small
\begin{tblr}{colspec={rcX[c]X[c]c},hlines=0pt,vlines=0pt,
  vline{3}={1}{text=\clap{\ce{+}}},
  vline{4}={1}{text=\clap{\ce{<=>[\qty{800}{\celsius}]}}},
  vline{5}={1}{text=\clap{\ce{+}}},
  }
  & \ce{CO} & \ce{H2O}（气） & \ce{CO2}& \ce{H2} \\
  反应开始时，各物质的浓度（\unit{mol/L}） & 0.01  & 0.01  & 0     & 0     \\
  达到平衡时，各物质的浓度（\unit{mol/L}） & 0.005 & 0.005 & 0.005 & 0.005 \\
\end{tblr}
}

为了进一步了解可逆反应达到化学平衡状态时的特点，对于上述反应，可以进行以下的实验。
设在四个 \qty{1}{L} 的密闭容器里，分别在每个容器里通入不同摩尔数的 \ce{CO}、\ce{H2O}（气）、\ce{CO2} 和 \ce{H2}，如\cref{tab:3-1} 中起始浓度一栏所表示的那样。
把四个密闭容器加热到 \qty{800}{\celsius}，经过足够的时间以后，如果这四个容器里的四种物质的浓度并不再随时间的改变而改变，那么，这四个容器里的反应混和物就都达到了平衡状态。
把实验测得的平衡浓度的数据列在\cref{tab:3-1} 的第二栏里。
为了研究在平衡混和物中各物质浓度之间的关系，我们用平衡混和物中各生成物的浓度乘积作为分子，把各反应物的浓度乘积作为分母，求其比值。并把求得的数值列入\cref{tab:3-1} 的第三栏里。

\begin{table}
  \caption{\ce{CO + H2O\,\text{（气）} <=> CO2 + H2 } 反应中起始和平衡时各物质的浓度（\qty{800}{\celsius}）}\label{tab:3-1}
  \begin{tblr}{colspec={*8{S[table-format=1.4,table-number-alignment = left]}X[c]},hline{3}=0.8pt,row{1,2}={m,c,guard}}
    \SetCell[c=4]{m,c}{起始时各物质的浓度（\unit{mol/L}）} & & & & \SetCell[c=4]{m,c}{平衡时各物质的浓度（\unit{mol/L}）}  & & & & \SetCell[r=2]{m,c}{ 平衡时\\[7pt] $\dfrac{[\ce{CO2}][\ce{H2}]}{[\ce{CO}][\ce{H2O}]}$} \\
    \ce{[CO]} & \ce{[H2O]} & \ce{[CO2]} & \ce{[H2]} & \ce{[CO]} & \ce{[H2O]} & \ce{[CO2]} & \ce{[H2]} & \\
    0.01   & 0.01 & 0      & 0      & 0.005  & 0.005  & 0.005  & 0.005  & 1.0 \\
    0      &    0 & 0.02   & 0.01   & 0.0067 & 0.0067 & 0.0133 & 0.0033 & 1.0 \\
    0.0025 & 0.03 & 0.0075 & 0.0075 & 0.0021 & 0.0296 & 0.0079 & 0.0079 & 1.0 \\
    0.01   & 0.03 & 0      & 0      & 0.0025 & 0.0225 & 0.0075 & 0.0079 & 1.0 \\
  \end{tblr}
\end{table}

从\cref{tab:3-1} 所列出的实验数据，可以得出如下的结论：

在一定温度下，可逆反应无论从正反应开始，或是从逆反应开始，又无论反应起始时反应物浓度的大小，最后都能达到平衡。
这时生成物的浓度乘积除以反应物的浓度乘积，便得到一个比值，而这个比值是个常数，这个常数叫做该反应的化学平衡常数（简称平衡常数）。
用符号 $K$ 表示。
例如，在 \qty{800}{\celsius} 时，
\[ \dfrac{[\ce{CO2}][\ce{H2}]}{[\ce{CO}][\ce{H2O}]}=1 \]
化学平衡常数 $K$ 的数值为 1。

\begin{Reading}[补充阅读]{}
为了更好地认识化学平衡状态的性质，下面我们还可以再研究一个可逆反应的例子。
\[ \ce{ H2 + I2\,\text{（气）} <=> 2HI }\]

在三个 \qty{1}{L} 的密闭容器里，分别在其中通入不同摩尔数的 \ce{H2}、\ce{I2} 或 \ce{HI}，如\cref{tab:3-2} 中起始浓度一栏所示。
把三个密闭容器都加热到 \qty{445}{\celsius}，经过足够的时间，使它们都达到平衡状态。
把实验测得的平衡状态时各物质的浓度数据列在\cref{tab:3-2} 的第二栏里。
然后，我们用平衡混和物中各生成物浓度与反应物浓度计算其不同的比值，并把求得的数值列入\cref{tab:3-2} 的第三栏和第四栏里。

\begin{tablehere}
  \begin{minipage}{\linewidth}\centering
  \caption{\ce{H2 + I2\,\text{（气）} <=> 2HI} 反应中起始和平衡时各物质的浓度（\qty{445}{\celsius}）}\label{tab:3-2}
  \begin{tblr}{colspec={*6{S[table-format=1.4,table-number-alignment = left]}X[c]X[c]},hline{3}=0.8pt,row{1,2}={m,c,guard}}
    \SetCell[c=3]{m,c}{ 起始时各物质的浓度 \\（\unit{mol/L}）}& & &\SetCell[c=3]{m,c}{ 起始时各物质的浓度 \\（\unit{mol/L}）} & & &\SetCell[r=2]{m,c}{平衡时\\[7pt] $\dfrac{[\ce{HI}]}{[\ce{H2}][\ce{I2}]}$}& \SetCell[r=2]{m,c}{平衡时\\[7pt] $\dfrac{[\ce{HI}]^2}{[\ce{H2}][\ce{I2}]}$} \\
    \ce{[H2]} & \ce{[I2]} & \ce{[HI]} & \ce{[H2]} & \ce{[I2]} & \ce{[HI]} & & \\
    0.01 & 0.01   & 0      & 0.0022 & 0.0022 & 0.0156 & 3323 & 50 \\
    0    & 0      & 0.01   & 0.0011 & 0.0011 & 0.0078 & 6446 & 50 \\
    0    & 0.0097 & 0.0349 & 0.0017 & 0.0114 & 0.0315 & 1625 & 51 \\
  \end{tblr}
\end{minipage}
\end{tablehere}

从\cref{tab:3-2} 所列出的数据可以看出，如果用平衡时反应物的浓度乘积去除生成物的浓度，如下式：
\[\frac{[\ce{HI}]}{[\ce{H2}][\ce{I2}]}\]
其比值就不是一个常数。但是如用反应物浓度的乘积 $[\ce{H2}][\ce{I2}]$ 去除生成物浓度的平方 $[\ce{HI}]^2$，即
\[\frac{[\ce{HI}]^2}{[\ce{H2}][\ce{I2}]}\]
我们发现在表里所列三种不同的情况下，仍然可以得到一个比例常数，即 $K=50$。这个常数就是这个反应在 \qty{445}{\celsius} 时的化学平衡常数。
\end{Reading}

对于一般的可逆反应
\[ \ce{ $m$A + $n$B <=> $p$C + $q$D }\]
其中 \ce{A}、\ce{B} 代表反应物，\ce{C}、\ce{D} 代表生成物，$m$、$n$、$p$、$q$ 分别表示化学方程式中各反应物和生成物的化学式前的系数。当在一定温度下达到平衡时，可以写出：
\[ \frac{[\ce{C}]^p[\ce{D}]^q}{[\ce{A}]^m[\ce{B}]^n}=K\]

$K$ 就叫做这个反应的平衡常数。

在平衡常数的表达式中，通常把生成物的浓度写在分式分子的位置，反应物的浓度写在分母的位置; 式中每种物质浓度的方次数就是化学方程式中该物质化学式前的系数。

从平衡常数 $K$ 值的大小,可以推断反应进行的程度。$K$ 值越大，表示化学反应达到平衡时生成物浓度对反应物浓度的比值越大，也就是反应进行得更完全些，反应物的转化率\footnote{$\text{某个指定反应物的转化率}=\dfrac{\text{指定反应物的起始浓度} - \text{指定反应物的平衡浓度}}{\text{指定反应物的起始浓度}}\times 100\%$}也越大。反之，$K$ 值越小，表示反应进行的程度越小，反应物的转化率也越小。

通过前面学习过的实验事实，我们知道在一定温度时，一个可逆反应的平衡常数并不随反应物或生成物浓度的改变而改变。但当温度改变时，情形就不同了，平衡常数 $K$ 值随温度的改变而改变。

例如，氢气跟碘化合生成碘化氢的反应：
\[ \ce{ H2\,\text{（气）} + I2\,\text{（气）} <=>[\triangle] 2HI\,\text{（气）} } \]
在不同温度时，其平衡常数不同，数值分别是：\qty{350}{\celsius}，$K=66.9$；\qty{425}{\celsius}，$K=54.4$；\qty{490}{\celsius}，$K=45.9$。

因此，在使用平衡常数时，必须注意是在哪一温度下进行的可逆反应。

\begin{example}
把 \qty{0.01}{mol} 的 \ce{CO2} 和 \qty{0.01}{mol} 的 \ce{H2} 放在 \qty{1}{L} 的密闭容器里，加热到 \qty{1200}{\celsius}，在生成 \qty{0.006}{mol} 的 \ce{CO} 和 \qty{0.006}{mol} 的 \ce{H2O}（气）的时候，反应达到平衡。试求这时该反应的化学平衡常数。
\end{example}
\begin{solution}\par
  \noindent{%\small
\begin{tblr}{colspec={rcX[c]X[c]c},hlines=0pt,vlines=0pt,
  vline{3}={1}{text=\clap{\ce{+}}},
  vline{4}={1}{text=\clap{\ce{<=>[\qty{1200}{\celsius}]}}},
  vline{5}={1}{text=\clap{\ce{+}}},
  }
  & \ce{CO} & \ce{H2O}（气）& \ce{CO2}& \ce{H2} \\
  反应开始时各物质的摩尔数 & 0.01  & 0.01  & 0     & 0     \\
  达到平衡时各物质的摩尔数 & $0.01-0.006$ & $0.01-0.006$ & 0.006 & 0.006 \\
\end{tblr}
}

  因在 \qty{1}{L} 的容器里，所以，平衡时各参加反应物质的浓度分别是：

  \ce{[CO]}=\qty{0.006}{mol/L}，\ce{[H2O]}=\qty{0.006}{mol/L}，\ce{[CO2]}=\qty{0.004}{mol/L}，\ce{[H2]}=\qty{0.004}{mol/L}，
  \[ K=\frac{[\ce{CO}][\ce{H2O}]}{[\ce{CO2}][\ce{H2}]}=\frac{0.006\times 0.006}{0.004\times 0.004}=2.25\]
  答: 在 \qty{1200}{\celsius} 时，这个反应的平衡常数是 2.25。
\end{solution}

\begin{example}
设在某温度时，在容积为 \qty{1}{L} 的密闭容器里，把 \ce{N2} 和 \ce{H2} 两种气体混和，反应后生成 \ce{NH3}。从实验测得，当达到平衡时，\ce{N2} 和 \ce{H2} 的浓度各为 \qty{2}{mol/L}，生成的 \ce{NH3} 的浓度是 \qty{3}{mol/L}，求这个反应在某温度的平衡常数和 \ce{N2}、 \ce{H2} 两气体在反应开始时的浓度。
\end{example}

\begin{solution}
  $\ce{N2 + 3H2 <=> 2NH3}$

  在某温度时，平衡混和物中 \ce{N2}、\ce{H2} 和 \ce{NH3} 的浓度分别是: \ce{[N2]}=\qty{2}{mol/L}，\ce{[H2]}=\qty{2}{mol/L}，\ce{[NH3]}=\qty{3}{mol/L}。
  \[K=\frac{[\ce{NH3}]^2}{[\ce{N2}][\ce{H2}]^3}=\frac{3^2}{2\times 2^3}=\frac{9}{16}=0.563\]

由化学方程式可以知道，生成 \qty{3}{mol} \ce{NH3}，必须消耗 \qty{1.5}{mol} 的 \ce{N2} 和 \qty{4.5}{mol} 的 \ce{H2}，\ce{N2} 和 \ce{H2} 在反应开始时的浓度应是消耗量和在平衡时剩余量之和。所以，
\begin{align*}
  \ce{N2}\,\text{在反应开始时的浓度}&=1.5+2=\qty{3.5}{mol/L} \\
  \ce{H2}\,\text{在反应开始时的浓度}&=4.5+2=\qty{6.5}{mol/L} 
\end{align*}
答: 这个反应在某温度的平衡常数是 0.563，\ce{N2} 在反应开始时的浓度为 \qty{3.5}{mol/L}，\ce{H2} 在反应开始时的浓度为 \qty{6.5}{mol/L}。
\end{solution}

\begin{example}
  在密闭容器中，给一氧化碳和水蒸气混和物加热时，达到下列平衡：
  \[ \ce{CO + H2O\,\text{（气）} <=> CO2 + H2}\]

  在 \qty{800}{\celsius} 时平衡常数等于 1，若用 \qty{2}{mol} 的 \ce{CO} 和 \qty{10}{mol} 的 \ce{H2O} 互相混和并加热到 \qty{800}{\celsius}，求 \ce{CO} 转化为 \ce{CO2} 的转化率？

  由于平衡常数表达式中各物质的浓度都是平衡时的浓度，因此在利用平衡常数的计算中，都必须首先列出平衡时各物质的浓度。
\end{example}
\begin{solution}
  设 $x=\text{达到平衡时 \ce{CO} 转化为 \ce{CO2} 的摩尔数}$，$V=\text{容器的容积}$

\noindent
\begin{tblr}{colspec={rcX[c]X[c]c},hlines=0pt,vlines=0pt,
  vline{3}={1}{text=\clap{\ce{+}}},
  vline{4}={1}{text=\clap{\ce{<=>}}},
  vline{5}={1}{text=\clap{\ce{+}}},
  }
  & \ce{CO} & \ce{H2O}（气）& \ce{CO2}& \ce{H2} \\
  反应开始时的浓度 & $\dfrac{2}{V}$   & $\dfrac{10}{V}$   & $\dfrac{0}{V}$ & $\dfrac{0}{V}$ \\
  达到平衡时的浓度 & $\dfrac{2-x}{V}$ & $\dfrac{10-x}{V}$ & $\dfrac{x}{V}$ & $\dfrac{x}{V}$ \\
\end{tblr}

  将平衡浓度代入平衡常数表达式，则得：
  \[K=\frac{[\ce{CO2}][\ce{H2}]}{[\ce{CO}][\ce{H2O}]}=\frac{\dfrac{x^2}{V^2}}{\dfrac{2-x}{V}\cdot\dfrac{10-x}{V}}=1\]
  \begin{align*}
    x^2&=(2-x)(10-x)\\
    x^2&=20-12x+x^2\\
    x&=\frac{20}{12}=1.66
  \end{align*}
  \ce{CO} 转化为 \ce{CO2} 的百分率为：
  \[ \frac{1.66}{2}\times 100\% =83\% \]
  答: \ce{CO} 转化为 \ce{CO2} 的转化率是 83\%。
\end{solution}

\begin{Practice}[习题]
\begin{question}
  \item 已知 $\ce{PCl5\,\text{（气）} <=> PCl3\,\text{（气）} + Cl2\,\text{（气）}}$
  在 \qty{230}{\celsius} 达到平衡时，平衡混和物中各物质的浓度分别是 \ce{[PCl5]}=\qty{0.47}{mol/L}，\ce{[PCl3]}=\qty{0.098}{mol/L}，\ce{[Cl2]}=\qty{0.098}{mol/L}，求这个反应的平衡常数。
  \item 已知 $\ce{H2\,\text{（气）} + I2\,\text{（气）} <=> 2HI\,\text{（气）}}$ 在 \qty{445}{\celsius} 时，平衡常数 $K=50$。计算在 \qty{445}{\celsius} 时，下面反应的平衡常数。
  \[\ce{ 2HI\,\text{（气）} <=> H2\,\text{（气）} + I2\,\text{（气）}}\]
  \item 在 \qty{25}{\celsius} 时，\ce{NO2} 与 \ce{N2O4} 处于平衡状态，两种气体的平衡混和物中含有 \qty{0.0125}{mol/L} 的 \ce{NO2} 和 \qty{0.0321}{mol/L}的 \ce{N2O4}。求这个反应的平衡常数。
  \item 在一定温度时，\qty{2}{mol} \ce{COCl2} 在 \qty{1}{L} 的密闭容器里分解，达到平衡时，测得已有 50\% 的 \ce{COCl2} 分解成 \ce{CO} 和 \ce{Cl2}。
  \[ \ce{ COCl2\,\text{（气）} <=> CO\,\text{（气）} + Cl2\,\text{（气）} }\]
  求在此温度时，这一反应的平衡常数。
  \item 在 \qty{1}{L} 的密闭容器里，装有 \qty{3}{mol} 的 \ce{NO2}，在一定温度时进行下面的反应: 
  \[\ce{2NO_2}\,\text{（气）} \ce{<=> N2O4}\,\text{（气）}\]
  反应的平衡常数 $K=7.15$，求在平衡时，该容器中的 \ce{NO2} 的摩尔数。
  \item 在密闭容器中将 \ce{NO2} 加热到某温度时，进行如下的反应：
  \[ \ce{ 2NO2 <=> 2NO +O2 } \]
  在平衡时各物质的浓度是: \ce{[NO2]}=\qty{0.06}{mol/L}，\ce{[NO]}=\qty{0.24}{mol/L}，\ce{[O2]}=\qty{0.12}{mol/L}。试求这温度下的平衡常数和 \ce{NO2} 在反应开始时的浓度。
  \item 已知一氧化碳跟水蒸气的反应为：
  \[ \ce{CO + H2O\,\text{（气）} <=> CO2 + H2 } \]
  在 \qty{427}{\celsius} 时的平衡常数为 9.4，若反应开始时，一氧化碳和水蒸气的浓度都是 \qty{0.01}{mol/L}，试计算在此反应条件下一氧化碳的转化率。
\end{question}
\end{Practice}

\section{影响化学平衡的条件}
化学平衡状态只有在一定的条件下才能保持。如果一个可逆反应达到平衡状态以后，反应条件（如浓度、压强、温度等）改变了，平衡混和物里各组成物质的百分含量也就随着改变而达到新的平衡状态，这叫做\Concept{化学平衡的移动}。

下面着重讨论浓度、压强和温度的改变对化学平衡的影响。

\subsection{浓度对化学平衡的影响}
当一个化学反应达到平衡的时候，其它反应条件不变，只改变其中任何一种反应物或生成物的浓度，就会改变正反应或逆反应的反应速度，使它们不再相等，从而使平衡移动。
\begin{Experiment}
  在一个小烧杯里混和 \qty{10}{ml} \qty{0.01}{M} 氯化铁溶液和 \qty{10}{ml} \qty{0.01}{M} 硫氰化钾溶液，溶液立即变成红色。

  把这红色溶液平均分到三个试管里，在第一个试管里加入少量 \qty{1}{M} 氯化铁溶液，在第二个试管里加入少量 \qty{1}{M} 硫氰化钾溶液。观察这两个试管里溶液颜色的变化，并跟第三个试管相比较。
\end{Experiment}

氯化铁跟硫氰化钾（\ce{KSCN}）起反应，生成红色的硫氰化铁 \ce{Fe(SCN)3}\footnote{实际上主要是 \ce{FeSCN^2+} 离子的颜色。}和氯化钾，这个反应可表示如下：
\[ \ce{ FeCl3 + 3KSCN <=> Fe(SCN)3 + 3KCl} \]

从上面实验可知，在平衡混和物里，当加入氯化铁溶液或硫氰化钾溶液以后，溶液的颜色都变深了。
这说明增大了任何一种反应物的浓度都促使化学平衡向正反应的方向移动，生成更多的硫氰化铁。

其它实验也可证明，在达到平衡的反应里，减小任何一种生成物的浓度，平衡会向正反应的方向移动; 减小任何一种反应物的浓度，平衡会向逆反应的方向移动。

由此可见，在其它条件不变的情况下，增大反应物的浓度，或减小生成物的浓度，都可以使平衡向着正反应的方向移动; 增大生成物的浓度或减小反应物的浓度，都可以使平衡向着逆反应的方向移动。

在生产上，往往采用增大容易取得的或成本较低的反应物浓度的方法，使成本较高的原料得到充分利用。
例如，在硫酸工业里，常用过量的空气使二氧化硫被充分氧化。

\subsection{压强对化学平衡的影响}
处于平衡状态的反应混和物里，不管是反应物或生成物，只要有气态物质存在，那么改变压强也常常会使化学平衡移动。
\begin{Experiment}\label{expr:3-5}
如\cref{fig:3-4}，用注射器（\qty{50}{ml} 或更大些的）吸入约 \qty{20}{ml} 二氧化氮和四氧化二氮的混和气体（使注射器的活塞达到 \MyRoman{1} 处）。吸入气体后，将细管端用橡皮塞加以封闭。然后把注射器的活塞往外拉到 \MyRoman{2} 处。观察当活塞反复地从 \MyRoman{2} 到 \MyRoman{1} 及从 \MyRoman{1} 到 \MyRoman{2} 时，管内混和气体颜色的变化。
\tcblower
\begin{figurehere}
  \includegraphics{3-4.pdf}
  \caption{压强对化学平衡的影响}\label{fig:3-4}
\end{figurehere}
\end{Experiment}

二氧化氮（棕色气体）跟四氧化二氮（无色气体）在一定条件下处于化学平衡状态。
在这个反应里，每减少二体积的 \ce{NO2}，就会生成一体积的 \ce{N2O4}。
\[ \ce{2NO2 \text{（气）} <=> N2O4 \text{（气）}} \]

从\cref{expr:3-5} 中可知，把注射器的活塞往外拉，管内体积增大，气体的压强减小，浓度减小，混和气体的颜色先变浅又逐渐变深。
逐渐变深是因为平衡向逆反应的方向移动，生成了更多的 \ce{NO2}。
把注射器的活塞往里压，管内体积减小，气体的压强增大，浓度增大，混和气体的颜色先变深又逐渐变浅。
逐渐变浅是因为平衡向正反应的方向移动，生成了更多的 \ce{N2O4}。

从上面的实验可以看出，在其它条件不变的情况下，增大压强会使化学平衡向着气体体积缩小的方向移动; 减小压强，会使平衡向着气体体积增大的方向移动。

在有些可逆反应里，反应前后气态物质的总体积没有变化，例如：
\[ \ce{2HI \text{（气）} <=> H2 \text{（气）} + I2 \text{（气）} } \]

在这种情况下，增大或减小压强就不能使化学平衡移动。

固态物质或液态物质的体积，受压强的影响很小，可以略去不计。
因此，平衡混和物都是固体或液体的，改变压强不能使平衡移动。

\subsection{温度对化学平衡的影响}
在吸热或放热的可逆反应里，反应混和物达到平衡状态以后，改变温度也会使化学平衡移动。
\begin{Experiment}
  装置如\cref{fig:3-5} 所示，两个连通着的烧瓶里，盛有二氧化氮跟四氧化二氮达到平衡的混和气体。然后用夹子夹住橡皮管把一个烧瓶放进热水里，把另一个烧瓶放入冰水（或冷水）里，观察混和气体的颜色的变化。
  \tcblower
  \begin{figurehere}
    \includegraphics{3-5.pdf}
    \caption{温度对化学平衡的影响}\label{fig:3-5}
  \end{figurehere}
\end{Experiment}

在二氧化氮生成四氧化二氮的反应里，正反应是放热反应，逆反应是吸热反应。
\[ \ce{ 2NO2 <=> N2O4 } + \qty{13.6}{kCal} \]

从上面实验可以知道，混和气体受热颜色变深，说明二氧化氮浓度增大，平衡向逆反应方向移动。
混和气体被冷却，颜色变浅，说明二氧化氮浓度减小，平衡向正反应的方向移动。

由此可见，在其他条件不变的情况下，温度升高，会使化学平衡向着吸热反应的方向移动; 温度降低，会使化学平衡向着放热反应的方向移动。

浓度、压强、温度对化学平衡的影响可以概括成一个原理来表示，这就是勒沙特列\footnote{勒沙特列（H.L. Le Chatelier, 1850--1936），法国化学家。}原理（亦称平衡移动原理）：\emph{如果改变影响平衡的一个条件（如浓度、压强或温度等），平衡就向能够减弱这种改变的方向移动}。

由于催化剂能够同样地增加正反应和逆反应的速度，因此它对化学平衡的移动没有影响，也就是说它不能改变达到化学平衡状态的反应混和物的百分组成。
但是使用了催化剂，就能够大大地缩短反应达到平衡所需的时间。

\begin{Practice}[习题]
\begin{question}
  \item 下述反应达到平衡时，
  \[ \ce{ 2NO + O2 <=> 2NO2} + \qty{27.02}{kCal}\]
  如果
  \begin{enumerate*}[(1)]
    \item 增大压强，
    \item 增大 \ce{O2} 的浓度，
    \item 减小 \ce{NO2} 的浓度，
    \item 升高温度，
  \end{enumerate*}
  平衡将向什么方向移动？并简述理由。
  \item 指出怎样改变反应物的浓度，可使下列反应的平衡向正反应方向移动？（\ce{C} 的浓度不变化）
  \[ \ce{ CO2 + C\,\text{（固）} <=> 2CO } \]

  如果升高温度，可使平衡向正反应方向移动，那么，生成一氧化碳的反应是放热反应还是吸热反应？
  \item 增大压强和升高温度对下面处于平衡状态的化学反应有什么影响？ 并说明原因。
  \begin{tasks}[label=(\arabic*)]
    \task \ce{ 2NO + O2 <=> 2NO2} ${}+\qty{27.02}{kCal}$ 
    \task \ce{ N2 + O2 <=> 2NO} ${}-\qty{43.2}{kCal}$
  \end{tasks}
  \item 化学反应
  \[ \ce{ CO\,\text{（气）} + NO2\,\text{（气）} <=> CO2\,\text{（气）} + NO\,\text{（气）} } + \qty{54.1}{kCal} \]
  在达到平衡时，如果
  \begin{enumerate*}[label=(\arabic*)] 
    \item 升高温度，
    \item 容器容积扩大到十倍，
  \end{enumerate*}
  平衡将分别受到什么影响？
  \item 在下面反应中，升高温度或降低温度对平衡有何影响？
  \[ \ce{ H2S <=> H2 + S\,\text{（固）} }-\qty{5}{kCal} \]
  \item 在密闭容器里，对一氧化碳和水蒸气的混和物加热时，达到下列平衡：
  \[ \ce{ CO + H2O\,\text{（气）} <=> CO2 + H2 }\]
  在 \qty{500}{\celsius} 时平衡常数 $K=9$，如果反应开始时，\ce{CO2} 及 \ce{H2O}（气）的浓度都是 \qty{0.02}{mol/L}，求 \ce{CO} 转化为 \ce{CO2} 的百分率？
  \item 在某温度下，当 \ce{ H2 + Br2\text{（气）} <=> 2HBr } 的反应达到平衡时，各反应物质的浓度分别是：
  \begin{align*}
    [\ce{H2}]&=\qty{0.5}{mol/L}\\
    [\ce{Br2}]&=\qty{0.1}{mol/L}\\
    [\ce{HBr}]&=\qty{1.6}{mol/L}
  \end{align*}
  求氢气和溴（气）的起始浓度和平衡常数。
\end{question}
\end{Practice}

\section{合成氨工业}
在这一节里，我们将讨论怎样应用学过的化学反应速度和化学平衡的知识来研究有关合成氨的一些问题。

\subsection{应用化学反应速度和化学平衡原理，选择合成氨的适宜条件}
氨的合成是一个放热的、气体总体积缩小的可逆反应。
\[\schemestart
\chemname{\ce{N2}}{\footnotesize （1 体积）} \+ 
\chemname{\ce{3H2}}{\footnotesize （3 体积）} 
\arrow(.mid east--.mid west){<=>}
\chemname{\ce{2NH3}}{\footnotesize （2 体积）} \+ \qty{22.08}{kCal}
\schemestop\chemnameinit{}\]

从勒沙特列原理来看，根据合成氨的反应式，降低温度、增大压强都会使平衡向着生成氨的方向移动，提高平衡混和物中氨的含量。

\cref{tab:3-3} 中的实验数据是在不同温度和压强下，平衡混和物中氨含量的变化情况。
\begin{table}
  \caption{达到平衡时混和气中氨的含量（体积百分数）（氮气和氢气的体积比是 $1:3$）}\label{tab:3-3}
  \begin{tblr}{colspec={c*6{X[r]}},row{1}={m,c},hline{2}=0.8pt}
    \diagboxthree{温度}{氨含量}{压强}&\qty{1}{atm} & \qty{100}{atm} & \qty{200}{atm} & \qty{300}{atm} & \qty{600}{atm} & \qty{1000}{atm}\\
    \qty{200}{\celsius}& 15.3 & 81.5 & 86.4 & 89.9 & 95.4 & 98.8 \\
    \qty{300}{\celsius}&  2.2 & 52.0 & 64.2 & 71.0 & 84.2 & 92.6 \\
    \qty{400}{\celsius}&  0.4 & 25.1 & 38.2 & 47.0 & 65.2 & 70.8 \\
    \qty{500}{\celsius}&  0.1 & 10.6 & 19.1 & 26.4 & 42.2 & 57.5 \\
    \qty{600}{\celsius}& 0.05 &  4.5 &  9.1 & 13.8 & 23.1 & 31.4 \\
  \end{tblr}
\end{table}

下面我们比较具体地研究改变压强、温度等条件以利提高合成氨的含量。
\begin{enumerate}[1. ]
  \item 压强  \quad 上面已讲过氨的合成是一个体积减小的反应，因此在温度一定时，增大压强有利于氨的合成。从\cref{tab:3-3} 中的实验数据也证实了这一点。但是压强越大，需要的动力越大，对材料的强度和设备的制造要求也越高。一般合成氨厂采用的压强是 \qtyrange{200}{500}{atm}。
  \item 温度  \quad 由于合成氨的反应是放热反应，因此当压强一定、温度升高时，氨的平衡浓度会降低，所以，从反应的理想条件来看，氨的合成反应，在较低温度下进行有利。\cref{tab:3-3} 中的实验数据也说明了这一点。但是温度过低，反应速度很慢，需要很长的时间才能达到平衡，这在工业上是很不经济的。所以,在实际生产中，合成氨反应是在 \qty{500}{\celsius} 左右的温度下进行的（所以选择 \qty{500}{\celsius} 左右，还因为工业上用的催化剂在这个温度活性最大）。
  \item 催化剂\quad 氮跟氢是极不容易化合的，即使在高压、高温下，虽可使反应速度加快一些，但仍然是十分缓慢的。通常，为了加快氮跟氢的合成反应速度，都采用加入催化剂以降低反应活化能的方法，使反应物在较低温度下，能较快地进行反应。
  
  目前，在工业上比较普遍地使用于合成氨的催化剂，主要是以铁为主体的多成分催化剂，又称铁触媒。
\end{enumerate}

在实际生产中，还需将生成的氨及时从混和气体中分离出来，并且不断地向循环气中补充氮气、氢气。

\subsection{合成氨工业简述}
合成氨工业主要是原料气中的氢气和氮气合成氨的过程。我们先从原料气谈起。

\subsubsection{原料气的制备、净化和压缩}
合成氨所需要的氮气，都取自空气。
从空气中制取氮气，通常有两种方法：一是将空气液化、蒸发，使它分离为氮和氧; 另一是将空气中的氧跟碳作用，生成二氧化碳，再除去二氧化碳，就可得到氮气。

氢来源于水和燃料。
可以将水蒸气通过赤热的煤（或焦炭）层，使水蒸气跟碳发生化学反应，生成一氧化碳和氢。
\[ \ce{C + H2O\,\text{（气）} \xlongequal{\triangle} CO + H2 } \]

工业生产中在催化剂作用下，一氧化碳还可以再跟水蒸气进行化学反应，生成二氧化碳和氢气。
除去二氧化碳后即可得到所需原料气。
\[ \ce{ CO + H2O\,\text{（气）}} \xlongequal[\triangle]{\text{催化剂}} \ce{CO + H2 } \]

石油、天然气、焦炉气（炼焦厂副产的气体）、炼厂气（石油炼制厂副产的气体）等燃料中都含有大量的碳氢化合物，这些碳氢化合物在一定条件下可以跟氧或者水蒸气进行化学反应，生成一氧化碳和氢。

在制取原料气的过程中，常混有不同数量的其他杂质，而合成氨所需要的是纯净的氮气和氢气，所以必须将这些杂质除去，否则这些杂质会使合成氨所用的催化剂 “中毒”。
清除杂质的过程叫做原料气的净化。

由于氨的合成是在高压下进行的，所以在氨的合成之前，还需要将氮、氢混和气体用压缩机压缩至高压。

如上所述可知，生产合成氨所需用的原料是空气、水和一些燃料（包括煤、天然气、石油等）。
在合成前需要进行原料气的制造、原料气的净化和压缩等过程。

\subsubsection{氨的合成}
\cref{fig:3-6} 是合成氨的简要流程示意图。
\begin{figure}
  \includegraphics{3-6.pdf}
  \caption{合成氨的简要流程示意图}\label{fig:3-6}
\end{figure}

图中，氮、氢混和气体经过压缩以后，进入氨合成塔。
化学反应在合成塔里进行，这个反应是在高压、适当的温度和有触媒存在的条件下进行的放热反应。

\begin{Reading}[补充阅读]{}
\par\medskip\noindent
\begin{minipage}{0.58\linewidth}\parindent2em
合成塔的构造除为了有利于催化，要能够安放厚层的触媒外，还要能够耐高压和能够调节温度。
为了符合这个要求，合成塔有耐高压的厚壁。
塔里主要有接触室和热交换器两部分。
氨的合成塔种类很多。
\cref{fig:3-7} 是一种合成塔的内部构造示意图。
合成塔是一个具有厚壁的钢筒，可以耐高压，塔壳外用绝缘材料包裹。
合成塔里上部为接触室，内放粒状触媒，下部是热交换器。

氮、氢混和气体是由塔的上部进入塔里，向下沿着塔壳和接触室间的间隙流到下部热交换器。
在热交换器里，混和气体在管子和管子之间的空隙曲曲折折地向上流动，同时，吸收管子里反应后的热气体的热量，这样混和气体的温度就逐渐升高。
预热了的氮、氢混和气体由塔的上部中间的粗管流进接触室。
然后由上而下地通过触媒层。
这时即有一部分氮气和氢气化合成氨。
反应后的气体经多孔板向下通入热交换器的管子，被管外新导入的混和气
\end{minipage}\hfill
\begin{minipage}{0.4\linewidth}
  \begin{figurehere}
  \includegraphics{3-7.pdf}
  \caption{一种合成塔的内部构造示意图}\label{fig:3-7}
\end{figurehere}
\end{minipage}\par\medskip\noindent
体所冷却，然后由气体出口离开合成塔。
在正常工艺操作的情况下，反应放出的热量足以保持接触室的温度，不需要另外供给热量。
如果反应物里的杂质已经仔细清除掉，触媒可以连续使用。
\end{Reading}

\subsubsection{氨的分离}
从合成塔里出来的混和气体，通常约含 15\% 左右的氨。
为了使氨从没有起反应的氮气和氢气里分离出来，要把混和气体通过冷凝器，使氨液化，然后在气体分离器里把液态氨分离出来，导入液氨贮罐。
由气体分离器出来的气体，经过循环压缩机，再送到合成塔里去。

使没有起反应的物质从反应后的生成物里分离出来，并重新回到反应器里去的工艺过程，叫做循环操作过程。
利用循环操作过程，即使平衡混和物里生成物的含量不很大，也可以充分地利用原料来制造所需的生成物。

我国合成氨工业的发展很迅速，它在支援农业生产、实现农业现代化中起着很重要的作用。
我国合成氨工业不但在产量上不断提高，而且在技术上也取得了比较显著的成就。

合成氨工业对化学工业和国防工业也都有重要的意义。

\begin{Reading}{化学模拟生物固氮简介}
我们已学习了合成氨的生产原理。
我们也知道氨和许多铵盐都是很重要的化学肥料。
这是因为氮是构成蛋白质的基本元素之一，是农作物生长的主要营养元素之一。
自然界里含有大量的氮元素，氮气占空气总体积的 78\%，但却不能被植物直接利用。
变成铵态的氮，就能被植物吸收。
但要把氢气和空气中的氮气转变成氨，却正如前面介绍的需要高压和一定的高温才行。
这就必须要有耐高压、高温的器材和设备，还需要大量动力。
那么能不能在常温常压条件下，把空气中的氮转变为铵态氮呢？
人们在六十多年来曾进行了大量的努力，希望在温和条件下实现氨的合成，但一直还没有成功。
然而，我们知道某些豆科植物，它们的根部有根瘤菌共生，根瘤菌能起固氮作用，即摄取空气中的氮转化成氨等，为植物直接吸收，这就叫做生物固氮现象，生物固氮是在常温、 常压下进行的。
实际上，地球上的氮气的固定，绝大部分是通过生物固氮进行的，据不完全统计，全世界工业合成氮肥中的氮只占固氮总量的 20\%。
那么，人们能不能向大自然学到这种本领呢？ 
这就需要研究如何模拟生物的功能，把生物的功能的原理用于化学，借以改善现有的并且创造崭新的化学工艺过程。
如果模拟成功，不仅可以大大提高氮肥工业的效率，发展农业生产，同时还会对很多化学工业产生深远影响。
\end{Reading}

\begin{Practice}[习题]
\begin{question}
  \item 在氮、氢合成氨的反应达到平衡时，反应混和物里各物质的浓度分别是: $[\ce{N2}]=\qty{3}{mol/L}$，$[\ce{H2}]=\qty{8}{mol/L}$，$[\ce{NH3}]=\qty{4}{mol/L}$。试求 \ce{N2}、\ce{H2} 的起始浓度和这个反应的平衡常数。
  \item 在制硫酸中，用空气来氧化 \ce{SO2} 生成 \ce{SO3}，
  \[ \ce{ 2SO2 + O2 <=>[\text{\qtyrange{400}{500}{\celsius}}][\ce{V2O5}] 2SO3 } + \qty{47}{kCal} \]
  为什么在生产上要用过量的空气，使用催化剂，并在适当的温度下进行？
\end{question}
\end{Practice}

\section*{内容提要}
\addcontentsline{toc}{section}{内容提要}\setcounter{subsection}{0}
化学反应速度是研究在单位时间内反应物或生成物浓度的变化，化学平衡是研究反应进行的方向和反应进行的程度。 
这两个问题对今后的化学学习和生产实践都是很重要的。

\subsection*{一、化学反应速度的表示方法}
化学反应速度通常是用单位时间内反应物浓度的减小或生成物浓度的增大来表示。
浓度一般用 \unit{mol/L} 表示; 时间可以根据反应的快慢用 \unit{s}、\unit{min} 或 \unit{h} 来表示。

\subsection*{二、影响化学反应速度的重要条件}
反应物能不能发生化学反应以及反应的快慢，首先决定于反应物质的性质。此外，化学反应速度还受反应进行时所处条件的影响，其中重要的是温度、浓度、压强和催化剂。

要反应物相互作用，必须使反应物分子相互碰撞，碰撞的分子要具有足够的能量，才能发生化学反应。
这种能够发生化学反应的分子碰撞，叫做有效碰撞。
活化分子具有的最低能量与分子平均能量的差叫做活化能。
反应物分子中活化分子的百分数越大，有效碰撞的次数越多，反应进行得就越快。
所以，一个反应进行的快慢，决定于有效碰撞的次数。
而增加反应的温度、浓度、压强（气体反应）和使用性能良好的催化剂，都能增加反应物分子间的有效碰撞次数，因而都能增大反应的速度。

\subsection*{三、化学平衡及化学平衡的移动}
在一定温度下，一个密闭容器里进行的可逆反应，当正反应和逆反应的反应速度相等时，达到化学平衡。
化学平衡是一种动态平衡。
在平衡时，反应混和物中各组成物质的浓度保持不变。

对于一般的可逆反应
\[ \ce{ $m$A + $n$B <=> $p$C + $q$D } \]
当达到平衡时，各反应物质的浓度满足下面的关系式：
\[\frac{[\ce{C}]^p[\ce{D}]^q}{[\ce{A}]^m[\ce{B}]^n}=K\]
$K$ 叫做这个反应的平衡常数。平衡常数随温度的不同而改变。

平衡移动原理（勒沙特列原理）：如果改变影响平衡的一个条件（如浓度、温度、压强等），平衡就向能够使这种改变减弱的方向移动。
还可具体分述如下：
\begin{itemize}
  \item 浓度
  \begin{itemize}
    \item 增大反应物浓度或减小生成物浓度——平衡向正反应方向移动
    \item 减小反应物浓度或增大生成物浓度——平衡向逆反应方向移动
  \end{itemize}
  \item 温度
  \begin{itemize}
    \item 升高温度——平衡向吸热方向移动
    \item 降低温度——平衡向放热方向移动
  \end{itemize}
  \item 压强
  \begin{itemize}
    \item 增大压强——平衡向气态物质体积缩小方向移动
    \item 减小压强——平衡向气态物质体积增大方向移动
  \end{itemize}
  \item 催化剂对化学平衡的移动没有影响。
\end{itemize}

\subsection*{四、合成氨反应及其适宜的条件}
由于用氮气和氢气合成氨的反应中，气体体积缩小和氮分子的不活动性，合成氨的工业生产过程需要在高压、适当温度（因系放热反应，温度不能过高）和有催化剂存在的条件下进行。

\begin{Review}
\begin{question}
  \item 某一化学反应
  \[ \ce{ A + B -> C}\]
  开始时，\ce{A} 的浓度为 \qty{2}{mol/L}，两分钟后反应达到平衡，\ce{A} 的浓度下降为 \qty{1.5}{mol/L}，求这个反应以 \ce{A} 的浓度变化来表示的反应速度是多少？
  \item 在 \ce{N2 + 3H2 <=> 2NH3} 反应列，经 \qty{2}{s} 后 \ce{NH3} 的浓度增大了 {0.6}{mol/L}，那么，在这 \qty{2}{s} 内，用 \ce{H2} 浓度变化表示的反应速度是多少？
  \item 把 \ce{HgCl2} 的晶体和 \ce{KI} 的晶体放在研钵中研磨，能缓慢地起反应，颜色逐渐变红。如果把 \ce{HgCl2} 溶液和 \ce{KI} 溶液混和，立即生成红色 \ce{HgI2} 沉淀。为什么？
  \item 下列数据是一些反应的平衡常数，试判断，哪个反应进行得最接近完全？哪个反应进行得最不完全？
  \begin{tasks}(3)
    \task $K=1$
    \task $K=10$
    \task $K=10^{-1}$
    \task $K=10^{10}$
    \task $K=10^{-10}$
  \end{tasks}
  \item 下列说法对不对？为什么？
  \begin{tasks}[label=(\arabic*)]
    \task 当可逆反应达到平衡时，逆反应就停止了。
    \task 在合成氨的反应中使用催化剂，可使氮的转化率提高。
    \task 温度升高，使吸热反应速度增大，使放热反应速度减小。
    \task 当温度不变时，增大反应物的浓度，使平衡常数减小; 增大生成物的浓度，使平衡常数增大。
    \task 增大压强对溶液中的酸碱中和反应没有什么影响。
  \end{tasks}
  \item 在某温度时，\ce{ H2 + I2 <=> 2HI} 的平衡常数是 50。在这温度下，使氢与碘蒸气发生反应。反应开始时，碘蒸气的浓度是 \qty{1}{mol/L}。当达到平衡时，碘化氢的浓度为 \qty{0.9}{mol/L}，求反应开始时氢的浓度是多少？
  \item 已知反应 \ce{  H2 \text{（气）} + I2 \text{（气）} <=> 2HI \text{（气）} } 在某温度达到平衡时，平衡常数 $K=62.5$，如果在一个 \qty{10}{L} 的容器里放入 \qty{5}{mol} \ce{H2} 和 \qty{5}{mol} \ce{I2}，于该温度达到平衡，计算平衡时 \ce{H2}、\ce{I2} 和 \ce{HI} 的浓度。
  \item \ce{ 2HI \text{（气）} <=> H2 \text{（气）} + I2 \text{（气）}  } 在 \qty{500}{\celsius} 时的 $K=0.016$。若反应开始时，在一个 \qty{2}{L} 容器内放入 \qty{1}{mol} \ce{HI}，求达到平衡时 \ce{H2}、\ce{I2} 和 \ce{HI} 的浓度。
  \item 在硫酸生产中，二氧化硫氧化生成三氧化硫的反应如下：
  \[ \ce{ 2SO2 +O2 <=> 2SO3} +\qty{47}{kCal} \]
  回答以下问题：
  \begin{tasks}[label=(\arabic*)]
    \task 升高温度对上述反应有什么影响？为什么工业上常在 \qtyrange{400}{500}{\celsius} 下进行反应？
    \task 增大压强对上述反应有什么影响？为什么工业上常在接近常压下操作？
    \task 从沸腾炉导出的炉气在进入接触室之前，为什么要进行除尘、洗涤、干燥等净化措施？
    \task 进入接触室的原料气中，为什么空气常是过量的？
    \task 为什么要用适当的催化剂如五氧化二钒？
    \task 从吸收塔放出的尾气中，为什么仍含有少量二氧化硫？
  \end{tasks}
\end{question}
\end{Review}
